\documentclass[12pt, titlepage]{article}

\usepackage{fullpage}
\usepackage[round]{natbib}
\usepackage{multirow}
\usepackage{booktabs}
\usepackage{tabularx}
\usepackage{graphicx}
\usepackage{float}
\usepackage{hyperref}
\hypersetup{
    colorlinks,
    citecolor=blue,
    filecolor=black,
    linkcolor=red,
    urlcolor=blue
}

\input{../../Comments.text}
\input{../../Common.text}

\newcounter{acnum}
\newcommand{\actheacnum}{AC\theacnum}
\newcommand{\acref}[1]{\hyperref[#1]{AC\ref*{#1}}}

\newcounter{ucnum}
\newcommand{\uctheucnum}{UC\theucnum}
\newcommand{\uref}[1]{\hyperref[#1]{UC\ref*{#1}}}

\newcounter{mnum}
\newcommand{\mthemnum}{M\themnum}
% \mref links to the module's entry in Section \ref{SecMD} (Module Decomposition)
\newcommand{\mref}[1]{\hyperref[#1:dec]{M\ref*{#1}}}

\begin{document}

\title{Module Guide for \progname{}} 
\author{\authname}
\date{\today}

\maketitle

\pagenumbering{roman}

\section{Revision History}

\begin{tabularx}{\textwidth}{p{3cm}p{2cm}X}
\toprule {\bf Date} & {\bf Version} & {\bf Notes}\\
\midrule
Date 1 & 1.0 & Notes\\
Date 2 & 1.1 & Notes\\
\bottomrule
\end{tabularx}

\newpage

\section{Reference Material}

This section records information for easy reference. 

\subsection{Abbreviations and Acronyms}

\renewcommand{\arraystretch}{1.2}
\begin{tabular}{l l} 
  \toprule		
  \textbf{symbol} & \textbf{description}\\
  \midrule 
  AC & Anticipated Change\\
  DAG & Directed Acyclic Graph \\
  M & Module \\
  MG & Module Guide \\
  OS & Operating System \\
  R & Requirement\\
  SC & Scientific Computing \\
  SRS & Software Requirements Specification\\
  \progname & Explanation of program name\\
  UC & Unlikely Change \\
  \wss{etc.} & \wss{...}\\
  \bottomrule
\end{tabular}\\

\newpage

\tableofcontents

\listoftables

\listoffigures

\newpage

\pagenumbering{arabic}

\section{Introduction}

Decomposing a system into modules is a commonly accepted approach to developing
software.  A module is a work assignment for a programmer or programming
team~\citep{ParnasEtAl1984}.  We advocate a decomposition
based on the principle of information hiding~\citep{Parnas1972a}.  This
principle supports design for change, because the ``secrets'' that each module
hides represent likely future changes.  Design for change is valuable for the
Troitsky Toolbox, which will be maintained by successive members of the club's
software subteam as competition rules change each year and new tools are added.

Our design follows the rules laid out by \citet{ParnasEtAl1984}, as follows:
\begin{itemize}
\item System details that are likely to change independently should be the
  secrets of separate modules.
\item Each data structure is implemented in only one module.
\item Any other program that requires information stored in a module's data
  structures must obtain it by calling access programs belonging to that module.
\end{itemize}

After completing the first stage of the design, the Software Requirements
Specification (SRS), the Module Guide (MG) is developed~\citep{ParnasEtAl1984}. The MG
specifies the modular structure of the system and is intended to allow both
designers and maintainers to easily identify the parts of the software.  The
potential readers of this document are as follows:

\begin{itemize}
\item New software subteam members: This document can be a guide for a new
  member of the club's software subteam taking over maintenance of the Toolbox
  to easily understand the overall structure and quickly find the relevant
  modules they are searching for.
\item Maintainers: The hierarchical structure of the module guide improves the
  maintainers' understanding when they need to make changes to the system. It is
  important for a maintainer to update the relevant sections of the document
  after changes have been made.
\item Designers: Once the module guide has been written, it can be used to
  check for consistency, feasibility, and flexibility. Designers can verify the
  system in various ways, such as consistency among modules, feasibility of the
  decomposition, and flexibility of the design.
\end{itemize}

The rest of the document is organized as follows.
Section \ref{SecChange} lists the anticipated and unlikely changes of the software requirements.
Section \ref{SecMH} summarizes the module decomposition that was constructed according to the anticipated changes.
Section \ref{SecConnection} specifies the connections between the software requirements and the modules.
Section \ref{SecMD} gives a detailed description of the modules.
Section \ref{SecTM} includes two traceability matrices. One checks the completeness of the design against the requirements provided in the SRS. The other shows the relation between anticipated changes and the modules.
Section \ref{SecUse} describes the use relation between modules. 
Section \ref{SecUI} presents the design of the user interfaces. 
Section \ref{SecComm} describes the design of the communication protocols.
Section \ref{SecDB} summarizes the database design.
Section \ref{SecTimeline} outlines the schedule of tasks and the team members responsible for them.
Finally, Section \ref{SecAI} discloses the use of generative AI tools in the preparation of this document.




\section{Anticipated and Unlikely Changes} \label{SecChange}
This section lists possible changes to the system. According to the likeliness of the change, the possible changes are classified into two categories. Anticipated changes are listed in Section \ref{SecAchange}, and unlikely changes are listed in Section \ref{SecUchange}.

\subsection{Anticipated Changes} \label{SecAchange}
Anticipated changes are the source of the information that is to be hidden inside the modules. Ideally, changing one of the anticipated changes will only require changing the one module that hides the associated decision. The approach adopted here is called design for change. Some anticipated changes affect more than one module; these cases are identified in the traceability matrix in Section \ref{SecTM}.

\begin{description}
\item[\refstepcounter{acnum} \actheacnum \label{acHardware}:]
  The hosting platform and deployment configuration of the system. This may change from a cloud service to a university server depending on the club's budget and the servers available to it.
\item[\refstepcounter{acnum} \actheacnum \label{acRules}:]
  The competition rule set. The rules are revised and republished before each year's competition.
\item[\refstepcounter{acnum} \actheacnum \label{acModelInput}:]
  The library used to read STEP files and the club conventions for how bridge models are exported. The library may be replaced if it cannot read all club exports, and the export conventions may be refined as teams adopt them.
\item[\refstepcounter{acnum} \actheacnum \label{acAnalysis}:]
  The structural analysis method and the set of performance metrics computed for a bridge. These will be iterated on to improve the accuracy of predictions, and new metrics will be added as the club identifies what is useful to teams.
\item[\refstepcounter{acnum} \actheacnum \label{acMaterial}:]
  The template for material test files and the method used to derive material properties from test results. These may change as the club adjusts its testing procedures or starts using new materials and suppliers.
\item[\refstepcounter{acnum} \actheacnum \label{acUI}:]
  The layout and design of the Toolbox Web App, including the library used to display 3D models. This will likely change based on user feedback to improve usability.
\item[\refstepcounter{acnum} \actheacnum \label{acServices}:]
  The set of tools offered by the Toolbox. Future software subteam members are expected to add new tools as the club's needs grow.
\item[\refstepcounter{acnum} \actheacnum \label{acDB}:]
  The database technology, schema, and backup strategy. The schema will grow as new tools and types of content are added, and the technology may change with the hosting platform.
\item[\refstepcounter{acnum} \actheacnum \label{acAuth}:]
  The authentication mechanism used to sign members in. This may be replaced by an external sign-in provider, such as McMaster single sign-on, to simplify account management.
\item[\refstepcounter{acnum} \actheacnum \label{acHandoff}:]
  The structure and templates of handoff documents. Executives are expected to refine these each year based on what incoming members found useful.
\item[\refstepcounter{acnum} \actheacnum \label{acSearch}:]
  The search strategy used to find content in the Knowledge Repository. This may need to change to keep searches fast as the amount of stored content grows.
\end{description}

\subsection{Unlikely Changes} \label{SecUchange}
The module design should be as general as possible. However, a general system is more complex. Sometimes this complexity is not necessary. Fixing some design decisions at the system architecture stage can simplify the software design. If these decisions should later need to be changed, then many parts of the design will potentially need to be modified. Hence, it is not intended that these decisions will be changed.

\begin{description}
\item[\refstepcounter{ucnum} \uctheucnum \label{ucIO}:]
  Input/Output devices (Input: keyboard, mouse, and file upload; Output: screen and downloaded files). Members interact with the system only through a web browser on a computer.
\item[\refstepcounter{ucnum} \uctheucnum \label{ucPlatform}:]
  Delivery as a web application for desktop browsers. Mobile devices are not supported, as the main tasks, such as uploading CAD models, are done on a computer.
\item[\refstepcounter{ucnum} \uctheucnum \label{ucBroker}:]
  The broker-based architecture, in which all communication between services passes through the Toolbox Broker. This is fundamental to letting tools be added without modifying existing services.
\item[\refstepcounter{ucnum} \uctheucnum \label{ucStack}:]
  The back-end technology stack of C\# on .NET. Changing it would require rewriting most of the system, and it was chosen to be maintainable by future subteam members.
\item[\refstepcounter{ucnum} \uctheucnum \label{ucModelFormat}:]
  STEP as the only accepted bridge model format. STEP is supported by all the modelling software the club's teams currently use.
\item[\refstepcounter{ucnum} \uctheucnum \label{ucCompetition}:]
  The overall format of the Troitsky competition. Individual rules change each year, but bridges are expected to remain built from the same materials and tested the same way.
\item[\refstepcounter{ucnum} \uctheucnum \label{ucRoles}:]
  The club roles on which access control is based. The club's organizational structure is not expected to change often.
\item[\refstepcounter{ucnum} \uctheucnum \label{ucUsers}:]
  Current club members as the only users of the system's services. Past members and the public have access only to the public landing page.
\end{description}



\section{Module Hierarchy} \label{SecMH}

This section provides an overview of the module design. Modules are summarized
in a hierarchy decomposed by secrets in Table \ref{TblMH}. The modules listed
below, which are leaves in the hierarchy tree, are the modules that will
actually be implemented.

\begin{description}
\item [\refstepcounter{mnum}\label{mHH}\mref{mHH}:] Hardware-Hiding Module
\item [\refstepcounter{mnum}\label{mWebApp}\mref{mWebApp}:] Web App Layout Module
\item [\refstepcounter{mnum}\label{mBroker}\mref{mBroker}:] Toolbox Broker Module
\item [\refstepcounter{mnum}\label{mAuthz}\mref{mAuthz}:] Authorization Rules Module
\item [\refstepcounter{mnum}\label{mMembers}\mref{mMembers}:] Member Registry Module
\item [\refstepcounter{mnum}\label{mContent}\mref{mContent}:] Repository Content Module
\item [\refstepcounter{mnum}\label{mModelImport}\mref{mModelImport}:] Model Import Module
\item [\refstepcounter{mnum}\label{mModelLibrary}\mref{mModelLibrary}:] Model Library Module
\item [\refstepcounter{mnum}\label{mMetrics}\mref{mMetrics}:] Performance Metrics Module
\item [\refstepcounter{mnum}\label{mRuleSet}\mref{mRuleSet}:] Rule Set Module
\item [\refstepcounter{mnum}\label{mCompliance}\mref{mCompliance}:] Compliance Check Module
\item [\refstepcounter{mnum}\label{mTestRecords}\mref{mTestRecords}:] Material Test Records Module
\item [\refstepcounter{mnum}\label{mMatProps}\mref{mMatProps}:] Material Properties Module
\item [\refstepcounter{mnum}\label{mHandoff}\mref{mHandoff}:] Handoff Documents Module
\item [\refstepcounter{mnum}\label{mAuthn}\mref{mAuthn}:] Authentication Module
\item [\refstepcounter{mnum}\label{mPersist}\mref{mPersist}:] Persistence Module
\item [\refstepcounter{mnum}\label{mSearch}\mref{mSearch}:] Search Module
\item [\refstepcounter{mnum}\label{mStep}\mref{mStep}:] STEP Parsing Module
\item [\refstepcounter{mnum}\label{mSolver}\mref{mSolver}:] Structural Solver Module
\item [\refstepcounter{mnum}\label{mViewer}\mref{mViewer}:] 3D Viewer Module
\end{description}


\begin{table}[h!]
\centering
\begin{tabular}{p{0.3\textwidth} p{0.6\textwidth}}
\toprule
\textbf{Level 1} & \textbf{Level 2}\\
\midrule


{Hardware-Hiding Module} & ~ \\
\midrule
\multirow{13}{0.3\textwidth}{Behaviour-Hiding Module} & Web App Layout Module\\
& Toolbox Broker Module\\
& Authorization Rules Module\\
& Member Registry Module\\
& Repository Content Module\\
& Model Import Module\\
& Model Library Module\\
& Performance Metrics Module\\
& Rule Set Module\\
& Compliance Check Module\\
& Material Test Records Module\\
& Material Properties Module\\
& Handoff Documents Module\\
\midrule
\multirow{6}{0.3\textwidth}{Software Decision Module} & Authentication Module\\
& Persistence Module\\
& Search Module\\
& STEP Parsing Module\\
& Structural Solver Module\\
& 3D Viewer Module\\
\bottomrule


\end{tabular}
\caption{Module Hierarchy}
\label{TblMH}
\end{table}

\section{Connection Between Requirements and Design} \label{SecConnection}

The design of the system is intended to satisfy the requirements developed in
the SRS. In this stage, the system is decomposed into modules. The connection
between requirements and modules is listed in Table~\ref{TblRT}.

\wss{The intention of this section is to document decisions that are made
  ``between'' the requirements and the design.  To satisfy some requirements,
  design decisions need to be made.  Rather than make these decisions implicit,
  they are explicitly recorded here.  For instance, if a program has security
  requirements, a specific design decision may be made to satisfy those
  requirements with a password.}

\section{Module Decomposition} \label{SecMD}

Modules are decomposed according to the principle of ``information hiding''
proposed by \citet{ParnasEtAl1984}. The \emph{Secrets} field in a module
decomposition is a brief statement of the design decision hidden by the
module. The \emph{Services} field specifies \emph{what} the module will do
without documenting \emph{how} to do it. For each module, a suggestion for the
implementing software is given under the \emph{Implemented By} title, listing
the languages, frameworks, and libraries used to implement the module. If the
entry is \emph{OS}, this means that the module is provided by the operating
system or by standard programming language libraries.

Only the leaf modules in the hierarchy have to be implemented. If a dash
(\emph{--}) is shown, this means that the module is not a leaf and will not have
to be implemented.

\subsection{Hardware Hiding Modules}

\subsubsection{Hardware-Hiding Module (M\ref*{mHH})}\label{mHH:dec}

\begin{description}
\item[Secrets:] The hosting platform and its deployment configuration, such as
  a cloud service (e.g., Azure App Service) or a university or club-owned server.
\item[Services:] Serves as the virtual hardware used by the rest of the system.
  This module provides the interface between the host machine and the software,
  supplying the runtime, file system, network connections, and HTTPS endpoint
  that the Toolbox is served from.
\item[Implemented By:] OS and hosting platform, packaged as a Docker container
\item[Type of Module:] Abstract Object
\end{description}

\subsection{Behaviour-Hiding Module}

\begin{description}
\item[Secrets:] The contents of the required behaviours.
\item[Services:] Includes programs that provide externally visible behaviour of
  the system as specified in the software requirements specification (SRS)
  documents. This module serves as a communication layer between the
  hardware-hiding module and the software decision module. The programs in this
  module will need to change if there are changes in the SRS.
\item[Implemented By:] --
\end{description}

\subsubsection{Web App Layout Module (M\ref*{mWebApp})}\label{mWebApp:dec}

\begin{description}
\item[Secrets:] The page layout, navigation structure, and visual design of the
  Toolbox Web App, and the way user input on a page is turned into a request.
\item[Services:] Serves the public landing page and the member pages for each
  tool. Accepts user input through forms and file uploads, passes each request
  to the Toolbox Broker with the signed-in member's identity, and displays the
  response or error message. Ends inactive sessions after the session timeout.
\item[Implemented By:] C\# with ASP.NET Core Razor Pages, HTML, and CSS
\item[Type of Module:] Abstract Object
\end{description}

\subsubsection{Toolbox Broker Module (M\ref*{mBroker})}\label{mBroker:dec}

\begin{description}
\item[Secrets:] The registry of available services and the protocol used to
  route requests to them, including timeouts and request logging.
\item[Services:] Lets each tool register its operations at startup. Checks every
  request against the Authorization Rules Module before forwarding it, directs
  authorized requests to the service that provides the operation, and returns
  the response. Returns an error when an operation does not exist or a service
  does not respond, so that the failure of one service does not affect the
  others. Logs every request it handles.
\item[Implemented By:] C\# on ASP.NET Core, exposing RESTful HTTP endpoints
\item[Type of Module:] Abstract Object
\end{description}

\subsubsection{Authorization Rules Module (M\ref*{mAuthz})}\label{mAuthz:dec}

\begin{description}
\item[Secrets:] The mapping from club roles and team assignments to the
  operations they are allowed to perform.
\item[Services:] Assigns exactly one role to each member account and decides
  whether a given role and team may perform a requested operation. Restricts a
  team's current-year models to the members of that team and the executives.
  Lets executives and maintainers change roles, with changes taking effect by the
  member's next request.
\item[Implemented By:] C\# using ASP.NET Core Identity roles and authorization policies
\item[Type of Module:] Abstract Object
\end{description}

\subsubsection{Member Registry Module (M\ref*{mMembers})}\label{mMembers:dec}

\begin{description}
\item[Secrets:] The data structure holding member accounts, their status, and
  their team assignment for each competition year.
\item[Services:] Provides access programs to add, update, deactivate, and look
  up registered members. Assigns each account to at most one team per year.
  Blocks deactivating an executive until a successor has been assigned to their
  position.
\item[Implemented By:] C\# on .NET, using ASP.NET Core Identity user accounts
\item[Type of Module:] Abstract Data Type
\end{description}

\subsubsection{Repository Content Module (M\ref*{mContent})}\label{mContent:dec}

\begin{description}
\item[Secrets:] The structure of a stored item and its metadata (title,
  category, year, team, and tags), and the rules for which metadata is
  required.
\item[Services:] Provides access to the Knowledge Repository. Accepts uploads of
  any file type and rejects those missing a title, a category from the fixed
  list, or a year. Lets members preview and download items. Lets the uploader or
  an executive edit or archive an item after confirmation, and excludes archived
  items from default searches. Reports an error when an item cannot be stored or
  retrieved.
\item[Implemented By:] C\# on .NET, using LINQ for data access
\item[Type of Module:] Abstract Data Type
\end{description}

\subsubsection{Model Import Module (M\ref*{mModelImport})}\label{mModelImport:dec}

\begin{description}
\item[Secrets:] The club conventions for how bridge models are exported from CAD
  software and how STEP files are accepted and returned.
\item[Services:] Lets members upload and download bridge models as STEP files
  (\texttt{.stp} or \texttt{.step}). Rejects files that are not valid STEP or
  contain no geometry, stops processing after the model analysis timeout, and
  explains why. Displays a summary of the imported model so the member can
  confirm it was read correctly.
\item[Implemented By:] C\# on ASP.NET Core, using its file upload and download handling
\item[Type of Module:] Abstract Object
\end{description}

\subsubsection{Model Library Module (M\ref*{mModelLibrary})}\label{mModelLibrary:dec}

\begin{description}
\item[Secrets:] The organization of bridge model records, including ownership by
  team, competition year, privacy, and the categorization of members and joints.
\item[Services:] Stores each uploaded model as its team's current model, visible
  only to that team. Lets members browse and load past designs read-only,
  filtered by year and team. Lets executives convert all current models into past
  designs at the end of the season.
\item[Implemented By:] C\# on .NET, using LINQ for data access
\item[Type of Module:] Abstract Data Type
\end{description}

\subsubsection{Performance Metrics Module (M\ref*{mMetrics})}\label{mMetrics:dec}

\begin{description}
\item[Secrets:] The set of performance metrics and the failure criterion used to
  judge a bridge, such as the factor of safety of each element.
\item[Services:] Estimates the failure load and failure location of a bridge and
  computes its load-to-mass ratio, using the material properties selected by
  year or batch. Displays the assumptions and limitations with each estimate,
  and reports an error instead of an estimate when the results are invalid.
  New metrics can be added without changing existing ones.
\item[Implemented By:] C\# on .NET
\item[Type of Module:] Library
\end{description}

\subsubsection{Rule Set Module (M\ref*{mRuleSet})}\label{mRuleSet:dec}

\begin{description}
\item[Secrets:] The data structure used to represent a competition rule set,
  where each rule is a property with an allowed value or range.
\item[Services:] Stores the rule set for each competition year. Lets executives
  import a year's rules and create, edit, and remove individual rules without
  changing code. Marks rules that cannot be checked from geometry as manual
  checks.
\item[Implemented By:] C\# on .NET, using LINQ for data access
\item[Type of Module:] Abstract Data Type
\end{description}

\subsubsection{Compliance Check Module (M\ref*{mCompliance})}\label{mCompliance:dec}

\begin{description}
\item[Secrets:] The algorithm used to evaluate a bridge model's properties
  against a rule set.
\item[Services:] Checks a bridge model against the active or a past year's rule
  set and reports pass or fail for each rule. For each failed rule, reports the
  rule, the measured value, and the allowed range. Displays the competition year
  of the rules used and lists the rules that must be checked manually.
\item[Implemented By:] C\# on .NET
\item[Type of Module:] Library
\end{description}

\subsubsection{Material Test Records Module (M\ref*{mTestRecords})}\label{mTestRecords:dec}

\begin{description}
\item[Secrets:] The format of material test records and the template used to
  import them.
\item[Services:] Records test results entered through a form or imported in bulk
  from CSV or Excel files. Converts values to base units, requires grain
  orientation for wood, and rejects records with missing fields or physically
  impossible values, identifying the rows at fault. Lets executives flag results
  as invalid without deleting them, compares results against previous years,
  and exports records to CSV.
\item[Implemented By:] C\# on .NET, with CSV and Excel (\texttt{.xlsx}) import and CSV export
\item[Type of Module:] Abstract Data Type
\end{description}

\subsubsection{Material Properties Module (M\ref*{mMatProps})}\label{mMatProps:dec}

\begin{description}
\item[Secrets:] The method used to derive material properties, such as elastic
  modulus and strength, from accepted test records.
\item[Services:] Aggregates the accepted test records for a material, year, or
  batch into a set of material properties and provides them to the Performance
  Metrics Module. Keeps a trace from each property back to the tests it was
  computed from.
\item[Implemented By:] C\# on .NET
\item[Type of Module:] Library
\end{description}

\subsubsection{Handoff Documents Module (M\ref*{mHandoff})}\label{mHandoff:dec}

\begin{description}
\item[Secrets:] The structure and templates of exit documents and role records.
\item[Services:] Provides an exit-document template for each executive position
  and team lead. Lets members fill out documents, save them as drafts, and submit
  them for review before they are published to the repository. Validates
  submitted documents against their template, maintains a role record for each
  executive position, shows executives which documents are still outstanding,
  and exports documents to PDF.
\item[Implemented By:] C\# with ASP.NET Core Razor Pages, with PDF export
\item[Type of Module:] Abstract Data Type
\end{description}

\subsection{Software Decision Module}

\begin{description}
\item[Secrets:] The design decision based on mathematical theorems, physical
  facts, or programming considerations. The secrets of this module are
  \emph{not} described in the SRS.
\item[Services:] Includes data structure and algorithms used in the system that
  do not provide direct interaction with the user. 
  % Changes in these modules are more likely to be motivated by a desire to
  % improve performance than by externally imposed changes.
\item[Implemented By:] --
\end{description}

\subsubsection{Authentication Module (M\ref*{mAuthn})}\label{mAuthn:dec}

\begin{description}
\item[Secrets:] The authentication mechanism, including how credentials are
  stored and verified and how role and membership data are retrieved from
  \mref{mAuthz} and \mref{mMembers}.
\item[Services:] Signs members in and out. Verifies credentials and rejects
  invalid or deactivated accounts, enforces the password policy, stores
  passwords only as salted hashes, and supplies the signed-in member's identity
  and role to the rest of the system.
\item[Implemented By:] C\# using ASP.NET Core Identity, with salted password hashing
\item[Type of Module:] Abstract Object
\end{description}

\subsubsection{Persistence Module (M\ref*{mPersist})}\label{mPersist:dec}

\begin{description}
\item[Secrets:] The database technology, schema, and backup strategy.
\item[Services:] Stores and retrieves the data of every other module, keeping
  records complete and consistent when several members use the system at once.
  Provides access programs for reading and writing records and stored files, and
  backs up the database at the backup interval.
\item[Implemented By:] Microsoft SQL Server, accessed from C\# through LINQ
\item[Type of Module:] Abstract Object
\end{description}

\subsubsection{Search Module (M\ref*{mSearch})}\label{mSearch:dec}

\begin{description}
\item[Secrets:] The query and indexing strategy used to find content in the
  database.
\item[Services:] Returns the stored items matching a keyword search, filtered by
  category, year, team, or tag, within the maximum search time.
\item[Implemented By:] C\# LINQ queries against Microsoft SQL Server
\item[Type of Module:] Library
\end{description}

\subsubsection{STEP Parsing Module (M\ref*{mStep})}\label{mStep:dec}

\begin{description}
\item[Secrets:] The library used to read STEP files and the method used to
  extract geometry and metadata from them.
\item[Services:] Reads a STEP file and provides its geometry and metadata in a
  structured form, including overall dimensions, mass, and the members and
  joints of the bridge, for use by the model, analysis, and compliance modules.
\item[Implemented By:] C\# using the IxMilia.Step library (fallback: Open CASCADE through C++/CLI)
\item[Type of Module:] Library
\end{description}

\subsubsection{Structural Solver Module (M\ref*{mSolver})}\label{mSolver:dec}

\begin{description}
\item[Secrets:] The structural analysis method, such as finite element analysis,
  and the library used to carry it out.
\item[Services:] Builds a structural model from a bridge's members, joints,
  material properties, supports, and applied load, and returns the resulting
  internal forces and stresses in each element for use by the Performance
  Metrics Module and the 3D Viewer Module.
\item[Implemented By:] C\# using the BriefFiniteElement.NET library (fallback:
  a direct stiffness solver built on Math.NET Numerics)
\item[Type of Module:] Library
\end{description}

\subsubsection{3D Viewer Module (M\ref*{mViewer})}\label{mViewer:dec}

\begin{description}
\item[Secrets:] The library and rendering method used to display 3D models in
  the browser.
\item[Services:] Displays a bridge model in an interactive 3D view that the
  member can rotate, zoom, and pan, with its units and a scale reference. Can
  highlight results from the Structural Solver Module, such as the predicted
  failure location, on the model.
\item[Implemented By:] JavaScript in the browser, using occt-import-js and three.js
\item[Type of Module:] Abstract Object
\end{description}

\section{Traceability Matrix} \label{SecTM}

This section shows two traceability matrices: between the modules and the
requirements and between the modules and the anticipated changes.

% the table should use mref, the requirements should be named, use something
% like fref
\begin{table}[H]
\centering
\begin{tabular}{p{0.2\textwidth} p{0.6\textwidth}}
\toprule
\textbf{Req.} & \textbf{Modules}\\
\midrule
R1 & \mref{mHH}, \mref{mInput}, \mref{mParams}, \mref{mControl}\\
R2 & \mref{mInput}, \mref{mParams}\\
R3 & \mref{mVerify}\\
R4 & \mref{mOutput}, \mref{mControl}\\
R5 & \mref{mOutput}, \mref{mODEs}, \mref{mControl}, \mref{mSeqDS}, \mref{mSolver}, \mref{mPlot}\\
R6 & \mref{mOutput}, \mref{mODEs}, \mref{mControl}, \mref{mSeqDS}, \mref{mSolver}, \mref{mPlot}\\
R7 & \mref{mOutput}, \mref{mEnergy}, \mref{mControl}, \mref{mSeqDS}, \mref{mPlot}\\
R8 & \mref{mOutput}, \mref{mEnergy}, \mref{mControl}, \mref{mSeqDS}, \mref{mPlot}\\
R9 & \mref{mVerifyOut}\\
R10 & \mref{mOutput}, \mref{mODEs}, \mref{mControl}\\
R11 & \mref{mOutput}, \mref{mODEs}, \mref{mEnergy}, \mref{mControl}\\
\bottomrule
\end{tabular}
\caption{Trace Between Requirements and Modules}
\label{TblRT}
\end{table}

\begin{table}[H]
\centering
\begin{tabular}{p{0.2\textwidth} p{0.6\textwidth}}
\toprule
\textbf{AC} & \textbf{Modules}\\
\midrule
\acref{acHardware} & \mref{mHH}\\
\acref{acInput} & \mref{mInput}\\
\acref{acParams} & \mref{mParams}\\
\acref{acVerify} & \mref{mVerify}\\
\acref{acOutput} & \mref{mOutput}\\
\acref{acVerifyOut} & \mref{mVerifyOut}\\
\acref{acODEs} & \mref{mODEs}\\
\acref{acEnergy} & \mref{mEnergy}\\
\acref{acControl} & \mref{mControl}\\
\acref{acSeqDS} & \mref{mSeqDS}\\
\acref{acSolver} & \mref{mSolver}\\
\acref{acPlot} & \mref{mPlot}\\
\bottomrule
\end{tabular}
\caption{Trace Between Anticipated Changes and Modules}
\label{TblACT}
\end{table}

% TODO (when writing this matrix): Section 4.1 states that anticipated changes
% affecting more than one module are identified here. The ACs were consolidated
% (e.g. acModelInput = STEP library + export conventions, acAnalysis = solver +
% metrics, acMaterial = test template + property derivation, acUI = layout + 3D
% viewer, acDB = database + schema + backups), so some will map to several
% modules. Add a short paragraph below Table TblACT listing each AC that maps to
% more than one module and why (MG checklist: one-to-one mapping, or explain the
% exceptions).
% TODO: The Introduction says Section \ref{SecTM} "includes two traceability
% matrices". If an Unlikely Changes -> Modules table is added (as in Becca's
% MG), update that sentence in the Introduction.

\section{Use Hierarchy Between Modules} \label{SecUse}

In this section, the uses hierarchy between modules is
provided. \citet{Parnas1978} said of two programs A and B that A {\em uses} B if
correct execution of B may be necessary for A to complete the task described in
its specification. That is, A {\em uses} B if there exist situations in which
the correct functioning of A depends upon the availability of a correct
implementation of B.  Figure \ref{FigUH} illustrates the use relation between
the modules. It can be seen that the graph is a directed acyclic graph
(DAG). Each level of the hierarchy offers a testable and usable subset of the
system, and modules in the higher level of the hierarchy are essentially simpler
because they use modules from the lower levels.

\wss{The uses relation is not a data flow diagram.  In the code there will often
be an import statement in module A when it directly uses module B.  Module B
provides the services that module A needs.  The code for module A needs to be
able to see these services (hence the import statement).  Since the uses
relation is transitive, there is a use relation without an import, but the
arrows in the diagram typically correspond to the presence of import statement.}

\wss{If module A uses module B, the arrow is directed from A to B.}

\begin{figure}[H]
\centering
%\includegraphics[width=0.7\textwidth]{UsesHierarchy.png}
\caption{Use hierarchy among modules}
\label{FigUH}
\end{figure}

%\section*{References}

\section{User Interfaces} \label{SecUI}

\wss{Design of user interface for software and hardware.  Attach an appendix if
needed. Drawings, Sketches, Figma}

\section{Design of Communication Protocols} \label{SecComm}

\wss{If there is communication between hardware and software, it may be necessary to explain the design of the communication protocol.  This section should be filled in with NA if it is not appropriate for a given project.}

N/A. The Troitsky Toolbox is a web application with no custom hardware
components, so there is no hardware-software communication protocol to design.


\section{Database Design} \label{SecDB}

\wss{If a database is a significant component of the project, the database
design should be summarized here.  This section should be filled in with NA if
it is not appropriate for a given project.}

\section{Timeline} \label{SecTimeline}

\wss{Schedule of tasks and who is responsible}

\wss{You can point to GitHub if this information is included there}

\section{Generative AI Use Disclosure} \label{SecAI}

\wss{ChatGPT (version 5) was used to brainstorm ideas for the template for this disclose section.}

\subsection{AI Tools Used}

\wss{Give the name of the AI tool(s) used and the version.}

\subsection{How and Where Used}

\wss{Explain what the tool(s) were used for.  Examples include: brainstorming,
explaining a technical concept, reviewing the document, generating or modifying
text, generating or modifying diagrams, identifying errors or omissions or
inconsistencies.}

\wss{For each of the uses, identify which parts of the document were affected.}

\wss{You do not need to report every interaction, but you should disclose the 
uses that materially contributed to your work}

\subsection{Verification of AI-Generated Content}

\wss{\begin{itemize}
    \item Describe how the team checked AI-generated or AI-assisted material
    before including it in the document.
    \item In particular, verify factual claims, technical assertions,
    requirements, calculations, references, and other information for which an
    AI system may produce incorrect or fabricated results.
    \item \textbf{The team is responsible for the accuracy, completeness,
consistency, and appropriateness of the submitted document, regardless of
whether material was generated by a human or an AI system.}
\end{itemize}}

\newpage{}

\bibliographystyle {plainnat}
\bibliography{../../../refs/References}

\end{document}